\section{Moving the scalar constant outside the norm tree}

This replacement starts from the frozen sixth-round construction with
builder SHA256
\texttt{3f63cb782d52857446cf39f6e545aeb74f380d1347826eafa888cbb487738d51}.
It removes one rational coordinate after the complete compiler while
preserving scalar order twenty-eight. The point charts, nonzero shifted-point
guard, scalar proxy, logarithms, detector tests, physical anchor, and scalar
pole equations are unchanged. The new scalar witnesses are justified by a
soundness and completeness transfer; a bijection between the old and new
scalar witness tuples is neither needed nor asserted. The inherited global
arithmetic remains conditional.

\subsection{The replacement and its quadratic implementation}

Retain the actual input polynomial and contact product
\[
 \beta=2401+2754(2X+1)^2
       =11016X^2+11016X+5155,
 \qquad r=\beta(\beta-1)(32\beta+343).
\]
For every rational input, $\beta\ge2401$ and $r>0$. Retain the scalar
points $H=[h]P$ and $Q=S-P=[a-1]P$, where $h\ne0$ and $a-1\ne0$,
and write
\[
 (A,B,C,D_0)=\bigl(f(H)_0,f(H)_1,f(Q)_0,f(Q)_1\bigr),
 \qquad T=\operatorname{Re}(k_Hw).
\]
These are the same functions as in the preceding construction:
$k_H=y(H)/x(H)$, $w=h_S\tau$, and $\tau^2=f(Q)$.
The existing inverse of $f(Q)_0$ still excludes $S=P$; the tangent case
$S=-P$ and all chart sign choices are retained.

Keep the three existing Laurent steps and scaled proxy,
\[
 rq_3=q_1^2,\quad rq_7=q_3^2,\quad rq_{15}=q_7^2,
 \quad q_1=r^{-1},\quad \kappa=q_{15}=r^{-15},
 \quad ry=Tq_{15},\quad y=T/r^{16}.
\]
Replace the scalar tree by
\begin{align*}
 n&=B^2/r^2+C^2/r,&
 L&=A^2/r+n^2/r^2,\\
 J&=D_0^2+\kappa y^2,&
 N&=L^2+\kappa J^2,\qquad W=1+N.
\end{align*}
The physical pole input remains $b=42rW^2$ and is passed through the
same normalized scalar pole equations. Since every coefficient in these
real norm expressions is positive, $N\ge0$ and $W\ge1$.

The two left nodes use exactly four rational coordinates:
\[
 z=B/r,\qquad rn=zB+C^2,\qquad u=n/r,\qquad rL=A^2+un.
\]
Each equation has degree at most two. The $J$ and $N$ nodes each retain
a stored square of their right argument, so their inverse-weight terms
are also genuinely quadratic in stored coordinates. The final addition
$W=1+N$ is affine and introduces no coordinate. Relative to the preceding
implementation, the former node $n_0=1+A^2/r$ disappears; the former
two-coordinate realization of $L=n_0^2+rn^2$ is replaced by the displayed
two-coordinate realization of the new $L$.

All new denominators are $r$, which is already positive and has an
explicit rational inverse. Thus the rational extensions defining the
new nodes are unique, and no new branch or zero-denominator domain is
introduced. The shifted-point identity guard is not inferred from
these scalar equations and is explicitly retained.

\subsection{Odd contact valuations are already positive}

Let $p$ be a consumed prime outside the retained exceptional set
$\mathcal S=\{2,3,5,7,13,17\}$. The three-contact argument supplies an
odd valuation $R=v_p(r)$. Its positivity need not be enforced by an
all-left constant path in a norm tree.

Indeed, if $v_p(X)<0$, the leading coefficient
$11016=2^3 3^4 17$ is a $p$-unit. The three terms in the expanded
polynomial for $\beta$ then have distinct relevant smallest valuation,
so
\[
 v_p(\beta)=2v_p(X)<0.
\]
Since $32$ is also a $p$-unit,
\[
 v_p(\beta-1)=v_p(32\beta+343)=2v_p(X),
 \qquad v_p(r)=6v_p(X),
\]
which is even. If instead $v_p(X)\ge0$, all three contact factors are
integral, so $v_p(r)\ge0$. Consequently oddness of $R$ implies $R>0$
for every consumed prime. This is a statement about the actual input
polynomial; it would not follow for an arbitrary independently chosen
rational $\beta$.

\subsection{Soundness with an outer constant}

At such a prime, $42$ is a unit and the unchanged scalar pole implication
gives
\[
 R+2v_p(W)>0,\qquad v_p(W)>-R/2.
\]
Because $R>0$, also $v_p(1)=0>-R/2$. Applying the ultrametric inequality
to $N=W-1$ therefore gives
\[
 v_p(N)>-R/2.
\]
This step allows cancellation between $1$ and $N$; no assertion of
noncancellation at the final affine addition is made.

Within every binary norm node, the two rational summands have opposite
valuation parity. In $n$ their coefficient valuations are $-2R$ and
$-R$; in $L$ they are $-R$ and $-2R$; in both $J$ and $N$ they are
$0$ and $-15R$. Thus each node's valuation is exactly the smaller
summand valuation. The argument also covers zero arguments by taking
$v_p(0)=+\infty$.

Descending from $N$ gives
\[
 v_p(L)>-R/4,\qquad v_p(J)>29R/4.
\]
The two left nodes give
\[
 v_p(A)>3R/8,\qquad v_p(n)>7R/8,\qquad
 v_p(B)>23R/16,\qquad v_p(C)>15R/16.
\]
The right node gives the same bounds as before,
\[
 v_p(D_0)>29R/8,\qquad v_p(y)>89R/8,
 \qquad v_p(T)>217R/8.
\]
In particular all four coefficients of $f(H)$ and $f(Q)$ have positive
integer valuation. Both points therefore belong to the formal group at
every place above $p$, before the scalar proxy lemma is applied.

The retained two-place proxy lemma, including its nonunit fixed-coordinate
case at $41$, now gives
\[
 v_p(T)=v_p\bigl((a-1)/h\bigr).
\]
Since $R\ge1$ and rational valuations are integral, the strict bound
$v_p(T)>217R/8$ implies $v_p(T)\ge28$. Since $h$ is a nonzero integer,
\[
 v_p(a-1)\ge28,
 \qquad M_{\mathcal S}(\beta)^{28}\mid a-1.
\]
This is the same uniform scalar order as before. No stronger bound
$v_p(T)\ge28R$ is asserted. In particular, the same scalar height
consumer remains valid without changing the logarithm precision or
introducing an exceptional prime.

\subsection{Completeness and the order of choices}

Fix an integral input, so $r$ is a fixed positive integer. The inherited
preparation chooses $H$ first and then chooses $Q$ sufficiently deeper
at a finite set of primes. For fixed $H$, the retained function
$T=\operatorname{Re}(k_Hw)$ tends to zero as $Q\to O$, because
$S\to P$, $h_S$ stays bounded, and $\tau\to0$. Arbitrarily large
finite lower bounds on the four scalar coefficients and on $T$ are
therefore available in the same order as before.

For $p\mid r$, put $R=v_p(r)\ge1$. The following explicit sufficient
bounds establish the new local preparation:
\[
 v_p(A)\ge R+1,\quad v_p(B),v_p(C)\ge2R+1,
 \quad v_p(D_0)\ge4R+1,\quad v_p(y)\ge12R+1.
\]
They imply
\[
 v_p(n)\ge2R+2,\qquad v_p(L)\ge R+2,
 \qquad v_p(J)\ge8R+2.
\]
Consequently
\[
 v_p(N)\ge\min(2R+4,R+4)=R+4>0,
 \qquad W\equiv1\pmod p.
\]
The lower bound for $y$ is just the finite requirement
$v_p(T)\ge28R+1$, attainable after $H$ is fixed. At a prime dividing
$42$ but not $r$, requiring the four coefficients and $y$ to be
divisible by $p$ again gives $N\equiv0$ and $W\equiv1$ modulo $p$.
These preparation statements use sufficiently deep formal neighborhoods;
they do not assume good reduction at the small or ramified primes.

Thus $42rW^2$ is positive and nonzero, its reduced denominator is a
square, and the same local squareclasses at $2,3,17$ and the same
normalized pole realization are available. The inequality $W\ge1$
retains the real preparation used previously. Once $S$ is fixed, every
changed node is a rational function of $H$ with nonzero denominators
at the prepared point. The finite lower bounds persist on an appropriate
congruence class of $h$. Roots and scales are therefore still fixed
before the contact-prime selection. No simultaneous prime selection
or new global theorem is used.

\subsection{Actual measured cost and finite checks}

With the sixth-round detector and anchor parameters unchanged, the
complete isolated pipeline gives
\[
 \text{raw }409/377,\qquad
 \text{generic compiler }399/367,\qquad
 \boxed{\text{final }395\text{ coordinates and }363\text{ residuals}}.
\]
The pairs list coordinates and residuals. The independent gate replay
checks all $348$ source gates and the compiler operations; four further
certified eliminations give the displayed final count. Every residual
is quadratic, so the sum of their squares remains quartic. These are
counts from the complete implementation, not estimates from the
displayed abbreviations.

The standalone geometry checker reconstructs the actual changed scalar
rows and their unchanged normalized pole consumer, as well as all
retained point, difference, logarithm, scalar proxy, and nonzero-domain
rows. It checks the exact input polynomial, affine valuation bounds,
and finite completeness inequalities. Its $223$ source assertions and
the retained exact identity, tangent, and split-$41$ checks supplement
this uniform proof; they do not establish the inherited global arithmetic.
